\subsection{散射矩阵与解析系数}
记 $\Sigma^\circ=\Sigma\mathbin{\backslash}\{0,\pm\ii q_0,z_a^\pm\}$. 在 $\Sigma^\circ$ 上, 两组 Jost 解为同一空间谱方程的基本解矩阵, 因而存在与 $x$ 无关的散射矩阵
\begin{equation}\label{eq:Sdef}
\phi_+(x,t,z)=\phi_-(x,t,z)S(z),\qquad S=(s_{jk})_{j,k=1}^{2}.
\end{equation}

\paragraph{为什么两组基底只相差一个常矩阵.}
先在非分支点处讨论. 下面将证明 $\det\phi_-=\gamma\ne0$,
所以散射关系等价于 $S=\phi_-^{-1}\phi_+$.
对恒等式 $\phi_-^{-1}\phi_-=I$ 求导,
\[
(\phi_-^{-1})_x\phi_-+\phi_-^{-1}(\phi_-)_x=0,
\qquad
(\phi_-^{-1})_x=-\phi_-^{-1}U.
\]
于是
\begin{align*}
S_x
&=(\phi_-^{-1})_x\phi_++\phi_-^{-1}(\phi_+)_x\\
&=-\phi_-^{-1}U\phi_++\phi_-^{-1}U\phi_+=0.
\end{align*}
这里 $U$ 可以随 $x$ 变化; 只因两组矩阵满足同一个空间方程,
这两个包含 $U$ 的项才精确相消.
把右乘矩阵逐元写出,
\[
(\phi_{-,1},\phi_{-,2})
\begin{pmatrix}s_{11}&s_{12}\\s_{21}&s_{22}\end{pmatrix}
=\bigl(s_{11}\phi_{-,1}+s_{21}\phi_{-,2},
 s_{12}\phi_{-,1}+s_{22}\phi_{-,2}\bigr).
\]
注意第一列的两个系数是 $s_{11},s_{21}$,
而不是散射矩阵的第一行. 所有后续反射系数的符号都依赖这一约定.

其列关系为
\begin{equation}\label{eq:scattering-columns}
\phi_{+,1}=s_{11}\phi_{-,1}+s_{21}\phi_{-,2},\qquad
\phi_{+,2}=s_{12}\phi_{-,1}+s_{22}\phi_{-,2}.
\end{equation}

\begin{proposition}[散射矩阵与 Wronskian]\label{prop:scattering}
在定理 \ref{thm:jost} 的条件下,
\begin{equation}\label{eq:detjost}
\det\phi_\pm=\det\mu_\pm=\gamma(z),\qquad \det S=1.
\end{equation}
散射系数满足
\begin{align}
s_{11}&=\dfrac{\det(\phi_{+,1},\phi_{-,2})}{\gamma},&
s_{22}&=\dfrac{\det(\phi_{-,1},\phi_{+,2})}{\gamma},\label{eq:wronskian}\\
s_{21}&=\dfrac{\det(\phi_{-,1},\phi_{+,1})}{\gamma},&
s_{12}&=\dfrac{\det(\phi_{+,2},\phi_{-,2})}{\gamma}.\nonumber
\end{align}
其中 $s_{11}$ 在 $D_-$ 解析, $s_{22}$ 在 $D_+$ 解析, 并连续延拓到相应闭域中避开 $0,\pm\ii q_0$ 的紧集. $s_{12},s_{21}$ 在 $\Sigma^\circ$ 上连续; 在共振原像处使用空间归一化值.
\end{proposition}
\begin{proof}

\par\medskip\noindent\textbf{Abel 恒等式在这个二阶系统中的完整求导.}\quad 
空间方程的四个分量是
\begin{align*}
(\phi_{11})_x&=-\ii k\phi_{11}+q\phi_{21},&
(\phi_{12})_x&=-\ii k\phi_{12}+q\phi_{22},\\
(\phi_{21})_x&=-\conj{q}\phi_{11}+\ii k\phi_{21},&
(\phi_{22})_x&=-\conj{q}\phi_{12}+\ii k\phi_{22}.
\end{align*}
这里 $\phi$ 可以逐次取 $\phi_+$ 或 $\phi_-$,
并未引入新的谱解. 对二阶行列式的两项分别求导:
\begin{align*}
\partial_x\det\phi
&=\partial_x(\phi_{11}\phi_{22}-\phi_{12}\phi_{21})\\
&=(-\ii k\phi_{11}+q\phi_{21})\phi_{22}
+\phi_{11}(-\conj{q}\phi_{12}+\ii k\phi_{22})\\
&\quad-(-\ii k\phi_{12}+q\phi_{22})\phi_{21}
-\phi_{12}(-\conj{q}\phi_{11}+\ii k\phi_{21})\\
&=-\ii k\phi_{11}\phi_{22}+\ii k\phi_{11}\phi_{22}
+\ii k\phi_{12}\phi_{21}-\ii k\phi_{12}\phi_{21}\\
&\quad+q\phi_{21}\phi_{22}-q\phi_{22}\phi_{21}
-\conj{q}\phi_{11}\phi_{12}+\conj{q}\phi_{12}\phi_{11}=0.
\end{align*}
每个矩阵元是标量, 所以最后两组也能交换次序而相消.
归一化端的行列式为
\begin{align*}
\det\phi_\pm
&=\det\mu_\pm\det e^{\ii\theta\sg}
=\det\mu_\pm\,e^{\ii\theta}e^{-\ii\theta}
=\det\mu_\pm,\\
\lim\limits_{x\to\pm\infty}\det\mu_\pm
&=\det\begin{pmatrix}1&-\dfrac{\ii q_0}{z}\\
-\dfrac{\ii q_0}{z}&1\end{pmatrix}
=1-\left(-\dfrac{\ii q_0}{z}\right)^2
=1+\dfrac{q_0^2}{z^2}=\gamma.
\end{align*}
由于行列式不随 $x$ 变化, 它等于自己的端点极限.
对散射关系取行列式,
\[
\gamma=\gamma\det S,
\qquad \gamma\ne0,
\qquad \det S=1.
\]

由 $\tr U=0$, Abel 恒等式给出 $\partial_x\det\phi_\pm=0$. 在归一化端取极限, 并使用 $\det e^{\ii\theta\sg}=1$, 得到 $\det\phi_\pm=\gamma$. 对 \eqref{eq:Sdef} 取行列式即得 $\det S=1$. 
\par\medskip\noindent\textbf{四个 Wronskian 分别怎样消去另一项.}\quad 
二列行列式对每列都是线性的, 且两列相同的行列式为零.
对第一条列关系, 两种取法为
\begin{align*}
\det(\phi_{+,1},\phi_{-,2})
&=s_{11}\det(\phi_{-,1},\phi_{-,2})
+s_{21}\det(\phi_{-,2},\phi_{-,2})\\
&=s_{11}\gamma+s_{21}\cdot0=s_{11}\gamma,\\
\det(\phi_{-,1},\phi_{+,1})
&=s_{11}\det(\phi_{-,1},\phi_{-,1})
+s_{21}\det(\phi_{-,1},\phi_{-,2})\\
&=s_{11}\cdot0+s_{21}\gamma=s_{21}\gamma.
\end{align*}
对第二条列关系, 应按下面的列次序取行列式:
\begin{align*}
\det(\phi_{-,1},\phi_{+,2})
&=s_{12}\det(\phi_{-,1},\phi_{-,1})
+s_{22}\det(\phi_{-,1},\phi_{-,2})\\
&=s_{22}\gamma,\\
\det(\phi_{+,2},\phi_{-,2})
&=s_{12}\det(\phi_{-,1},\phi_{-,2})
+s_{22}\det(\phi_{-,2},\phi_{-,2})\\
&=s_{12}\gamma.
\end{align*}
例如, 若把最后一个行列式的两列反过来,
会得到 $-s_{12}\gamma$, 因而不能随意交换列序.
除以 $\gamma$ 后恰好得到原文四个散射系数.

对第一条列关系分别与 $\phi_{-,2}$、$\phi_{-,1}$ 取行列式, 有
\[
\det(\phi_{+,1},\phi_{-,2})=s_{11}\gamma,\qquad
\det(\phi_{-,1},\phi_{+,1})=s_{21}\gamma.
\]
第二条列关系同理给出 $s_{12},s_{22}$, 即 \eqref{eq:wronskian}. 
\par\medskip\noindent\textbf{把四个行列式中的指数也逐个抽出.}\quad 
由第一列与第二列分别带 $e^{\ii\theta}$ 与 $e^{-\ii\theta}$,
\begin{align*}
\det(\phi_{+,1},\phi_{-,2})
&=e^{\ii\theta}e^{-\ii\theta}
\det(\mu_{+,1},\mu_{-,2})
=\det(\mu_{+,1},\mu_{-,2}),\\
\det(\phi_{-,1},\phi_{+,2})
&=e^{\ii\theta}e^{-\ii\theta}
\det(\mu_{-,1},\mu_{+,2})
=\det(\mu_{-,1},\mu_{+,2}),\\
\det(\phi_{-,1},\phi_{+,1})
&=e^{2\ii\theta}\det(\mu_{-,1},\mu_{+,1}),\\
\det(\phi_{+,2},\phi_{-,2})
&=e^{-2\ii\theta}\det(\mu_{+,2},\mu_{-,2}).
\end{align*}
因此四个系数也可写为
\begin{align*}
s_{11}&=\dfrac{(\mu_+)_{11}(\mu_-)_{22}
-(\mu_-)_{12}(\mu_+)_{21}}{\gamma},\\
s_{22}&=\dfrac{(\mu_-)_{11}(\mu_+)_{22}
-(\mu_+)_{12}(\mu_-)_{21}}{\gamma},\\
s_{21}&=\dfrac{e^{2\ii\theta}}{\gamma}
\bigl((\mu_-)_{11}(\mu_+)_{21}
-(\mu_+)_{11}(\mu_-)_{21}\bigr),\\
s_{12}&=\dfrac{e^{-2\ii\theta}}{\gamma}
\bigl((\mu_+)_{12}(\mu_-)_{22}
-(\mu_-)_{12}(\mu_+)_{22}\bigr).
\end{align*}
$s_{22}$ 的两列恰为 $\mu_{-,1},\mu_{+,2}$,
二者都在 $D_+$ 解析; $s_{11}$ 的两列恰为
$\mu_{+,1},\mu_{-,2}$, 二者都在 $D_-$ 解析.
$\gamma$ 在相应开域没有零点, 所以除法保留解析性.
另一方面, $s_{21}$ 与 $s_{12}$ 所用两列来自相反解析域,
这里只能在共同边界上取值, 不能由这两个式子直接宣称全域解析.

对透射系数, 两列的指数相消, 故
\[
s_{11}=\dfrac{\det(\mu_{+,1},\mu_{-,2})}{\gamma},\qquad
s_{22}=\dfrac{\det(\mu_{-,1},\mu_{+,2})}{\gamma}.
\]
两个分子分别由 $D_-$ 和 $D_+$ 的解析列组成, 且 $\gamma$ 的零点只位于分支点. 定理 \ref{thm:jost} 因而给出所述解析性与连续性. 对非对角元, 列指数只留下连续因子 $e^{\pm2\ii\theta}$, 故它们在 $\Sigma^\circ$ 上连续.
\end{proof}

按照 $\phi_+=\phi_-S$ 的约定, 以下称 $s_{11},s_{22}$ 为透射系数. 在其非零的谱弧上定义反射系数
\begin{equation}\label{eq:reflection}
\rho(z)=\dfrac{s_{21}(z)}{s_{11}(z)},\qquad
\widetilde\rho(z)=\dfrac{s_{12}(z)}{s_{22}(z)}.
\end{equation}
由 $\det S=1$,
\begin{equation}\label{eq:scattering-determinant}
s_{11}s_{22}-s_{12}s_{21}=1,\qquad
s_{11}s_{22}(1-\rho\widetilde\rho)=1.
\end{equation}

\paragraph{反射系数与行列式约束的代数.}
由定义, 在分母不为零的谱弧上有
\[
s_{21}=\rho s_{11},\qquad s_{12}=\widetilde\rho s_{22}.
\]
将这两式原次序代入行列式,
\begin{align*}
s_{11}s_{22}-s_{12}s_{21}
&=s_{11}s_{22}-(\widetilde\rho s_{22})(\rho s_{11})\\
&=s_{11}s_{22}(1-\rho\widetilde\rho)=1.
\end{align*}
由于这里都是标量函数, 可以交换乘法次序.
这也说明 $1-\rho\widetilde\rho$ 不能在该谱弧上为零.
共轭约化之后, 实轴上会有
$\widetilde\rho=-\conj\rho$, 从而该因子变成
$1+|\rho|^2$; 这里先保留原有符号形式,
以免把只在实轴成立的等式用于圆谱的非实点.

后续矩阵 RHP 取普通连续谱上 $s_{11}s_{22}\ne0$ 的散射数据; 分支点的边界正则性另行列入 RHP 的条件.

\subsection{演化与共振相容性}
设三场为包含 $t=0$ 的开区间 $I$ 上的经典规范 HMB 解, 并满足
\begin{equation}\label{eq:evolution-regularity}
\begin{gathered}
q-q_0\in C(I;W^{3,1}(\mathbb R))\cap C^1(I;L^1(\mathbb R)),\\
p+\dfrac{\ii\eta_0q_0}{a},\quad\eta-\eta_0\in C(I;W^{1,1}(\mathbb R)).
\end{gathered}
\end{equation}
这些条件保证可在避开共振原像的谱参数紧集上对 Jost 积分方程求 $t$ 导数, 且 $V-V_\infty$ 在空间两端趋于零.

\begin{proposition}[散射数据的时间演化]\label{prop:evolution}
在条件 \eqref{eq:evolution-regularity} 下, 共振点之外的 $\phi_\pm$ 同时满足 Lax 对 \eqref{eq:gaugelax}, $S$ 与 $x,t$ 无关. 对固定 $t$ 的空间归一化 Jost 解
\[
\chi_\pm=\phi_\pm e^{-\ii\lambda Ft\sg},\qquad
\chi_\pm e^{\ii\lambda x\sg}\longrightarrow Y\quad(x\to\pm\infty),
\]
定义 $\chi_+=\chi_-S_{\rm sp}$. 则
\begin{equation}\label{eq:timeS}
S_{\rm sp}(z,t)=e^{\ii\lambda Ft\sg}S(z)e^{-\ii\lambda Ft\sg},
\end{equation}
\begin{equation}\label{eq:timer}
\rho_{\rm sp}(z,t)=e^{-2\ii\lambda Ft}\rho_{\rm sp}(z,0),\qquad
\widetilde\rho_{\rm sp}(z,t)=e^{2\ii\lambda Ft}\widetilde\rho_{\rm sp}(z,0).
\end{equation}
透射系数、离散零点及同时归一化的离散列相关常数均与 $t$ 无关.
\end{proposition}
\begin{proof}

\par\medskip\noindent\textbf{零曲率怎样保证时间偏差仍是空间解.}\quad 
$\chi_\pm$ 满足 $(\chi_\pm)_x=U\chi_\pm$.
交换混合导数并逐项求导,
\begin{align*}
\partial_x(\chi_{\pm,t}-V\chi_\pm)
&=\partial_t(U\chi_\pm)-V_x\chi_\pm-VU\chi_\pm\\
&=U\chi_{\pm,t}+(U_t-V_x-VU)\chi_\pm\\
&=U\chi_{\pm,t}-UV\chi_\pm\\
&=U(\chi_{\pm,t}-V\chi_\pm).
\end{align*}
第三行使用
$U_t-V_x+[U,V]=U_t-V_x+UV-VU=0$.
空间归一化中没有 $t$ 指数, 因而在相应端点
\[
\chi_\pm\sim Ye^{-\ii\lambda x\sg},
\qquad
\chi_{\pm,t}\sim0.
\]
另一方面,
\begin{align*}
V_\infty Ye^{-\ii\lambda x\sg}
&=\ii\lambda F Y\sg e^{-\ii\lambda x\sg}\\
&=Ye^{-\ii\lambda x\sg}\ii\lambda F\sg.
\end{align*}
所以 $\chi_{\pm,t}-V\chi_\pm$ 的端态恰好等于
$-\chi_\pm\ii\lambda F\sg$ 的端态.
两者均满足同一个空间方程, 用 Jost 唯一性比较每一列即可相等.
正则性条件的用途是保证先前交换 $x,t$ 微分和求端点极限合法;
不能仅凭形式上的零曲率就跳过这些步骤.

由零曲率条件,
\[
\partial_x(\chi_{\pm,t}-V\chi_\pm)=U(\chi_{\pm,t}-V\chi_\pm).
\]
在普通连续谱上, $V_\infty Y=\ii\lambda F Y\sg$. 由 \eqref{eq:evolution-regularity} 和归一化的 $t$ 导数, 空间解 $\chi_{\pm,t}-V\chi_\pm$ 在归一化端与 $-\chi_\pm\ii\lambda F\sg$ 具有相同渐近行为. 对相应齐次 Volterra 方程使用唯一性, 得到
\[
\chi_{\pm,t}=V\chi_\pm-\chi_\pm\ii\lambda F\sg.
\]

\par\medskip\noindent\textbf{恢复共同传播相位时的乘积求导.}\quad 
因为 $\lambda,F$ 都与 $t$ 无关,
\begin{align*}
\partial_t\bigl(\chi_\pm e^{\ii\lambda Ft\sg}\bigr)
&=\chi_{\pm,t}e^{\ii\lambda Ft\sg}
+\chi_\pm\ii\lambda F\sg e^{\ii\lambda Ft\sg}\\
&=(V\chi_\pm-\chi_\pm\ii\lambda F\sg)e^{\ii\lambda Ft\sg}
+\chi_\pm\ii\lambda F\sg e^{\ii\lambda Ft\sg}\\
&=V\chi_\pm e^{\ii\lambda Ft\sg}=V\phi_\pm.
\end{align*}
这两个 $\ii\lambda F\sg$ 项的符号必须相反才会抵消.
空间关系也保持不变, 因为右侧指数不依赖 $x$.

右乘 $e^{\ii\lambda Ft\sg}$, 得到 $\phi_{\pm,t}=V\phi_\pm$. 同一论证作用于各解析列. 对 \eqref{eq:Sdef} 求导, 有
\[
\phi_-S_x=U\phi_+-U\phi_-S=0,\qquad
\phi_-S_t=V\phi_+-V\phi_-S=0.
\]

\par\medskip\noindent\textbf{从矩阵恒定性到逐元的时间演化.}\quad 
$\det\phi_-\ne0$, 所以原文两条矩阵等式左乘
$\phi_-^{-1}$ 后给出 $S_x=S_t=0$.
将 $\phi_\pm=\chi_\pm e^{\ii\lambda Ft\sg}$
代回散射关系, 得
\begin{align*}
\chi_+e^{\ii\lambda Ft\sg}
&=\chi_-e^{\ii\lambda Ft\sg}S,\\
\chi_+&=\chi_-e^{\ii\lambda Ft\sg}Se^{-\ii\lambda Ft\sg}.
\end{align*}
因此 $S_x=S_t=0$. 由 $\chi_\pm=\phi_\pm e^{-\ii\lambda Ft\sg}$ 及 $\phi_+=\phi_-S$,
\[
\chi_+=\chi_-e^{\ii\lambda Ft\sg}S e^{-\ii\lambda Ft\sg},
\]
与 $\chi_+=\chi_-S_{\rm sp}$ 比较即得 \eqref{eq:timeS}. 展开右端,
\begin{align*}
S_{\rm sp}
&=\begin{pmatrix}e^{\ii\lambda Ft}&0\\0&e^{-\ii\lambda Ft}\end{pmatrix}
\begin{pmatrix}s_{11}&s_{12}\\s_{21}&s_{22}\end{pmatrix}
\begin{pmatrix}e^{-\ii\lambda Ft}&0\\0&e^{\ii\lambda Ft}\end{pmatrix}\\
&=\begin{pmatrix}
s_{11}&e^{2\ii\lambda Ft}s_{12}\\
e^{-2\ii\lambda Ft}s_{21}&s_{22}
\end{pmatrix}.
\end{align*}
由对角元可见, 空间归一化的透射系数仍为 $s_{11},s_{22}$, 因而它们的解析延拓与离散零点均不随 $t$ 改变. 再由非对角元及反射系数的定义 \eqref{eq:reflection}, 并使用 $S_{\rm sp}(z,0)=S(z)$, 得到 \eqref{eq:timer}. 同时归一化的离散列相关常数不随 $t$ 改变, 则由两列共同满足的传播方程得到.

所以两个对角元恒定, 非对角元分别乘相反符号的指数.
按照原文的比值定义再除以不变的对角元,
\begin{align*}
\rho_{\rm sp}(t)
&=\dfrac{e^{-2\ii\lambda Ft}s_{21}}{s_{11}}
=e^{-2\ii\lambda Ft}\rho_{\rm sp}(0),\\
\widetilde\rho_{\rm sp}(t)
&=\dfrac{e^{2\ii\lambda Ft}s_{12}}{s_{22}}
=e^{2\ii\lambda Ft}\widetilde\rho_{\rm sp}(0).
\end{align*}
透射系数作为解析函数恒定, 因而它的零点与对 $z$ 的导数都恒定.
离散列的相关常数所指的是同时满足两条 Lax 方程的 $\phi$ 列;
若改用只作空间归一化的 $\chi$ 列, 相关常数会含传播指数.
\end{proof}

共振点的空间归一化取值由 \eqref{eq:volterra} 定义, 介质方程在此固定了散射矩阵的形式.

\begin{proposition}[共振点的反射条件]\label{prop:resonance}
固定 $t$, 设 $q-q_0\in L^1(\mathbb R)$, 三场满足介质方程及共同端态
\[
(q,p,\eta)\longrightarrow\left(q_0,-\dfrac{\ii\eta_0q_0}{a},\eta_0\right),\qquad x\to\pm\infty.
\]
若 $z_a^\pm$ 处的空间 Jost 基底 $\chi_\pm$ 连续可逆, 且具有上述振荡归一化, 则
\[
[S_{\rm sp}(z_a^\pm,t),\sg]=0,\qquad
\rho_{\rm sp}(z_a^\pm,t)=\widetilde\rho_{\rm sp}(z_a^\pm,t)=0.
\]
\end{proposition}
\begin{proof}

\par\medskip\noindent\textbf{介质约束在共振谱点如何变成守恒矩阵.}\quad 
这里始终固定 $k=a$, 并且只使用空间解 $\chi_\pm$,
因此没有计算在 $F$ 的极点上的指数.
逆矩阵的空间导数是
\[
(\chi_\pm^{-1})_x=-\chi_\pm^{-1}U(a).
\]
对被共轭的介质矩阵逐项求导,
\begin{align*}
\partial_x(\chi_\pm^{-1}R\chi_\pm)
&=-\chi_\pm^{-1}U(a)R\chi_\pm
+\chi_\pm^{-1}R_x\chi_\pm
+\chi_\pm^{-1}RU(a)\chi_\pm\\
&=\chi_\pm^{-1}\bigl(R_x-U(a)R+RU(a)\bigr)\chi_\pm\\
&=\chi_\pm^{-1}\bigl(R_x-[U(a),R]\bigr)\chi_\pm=0.
\end{align*}
这说明它不随 $x$ 变化, 可在任一归一化端计算.

由 $R_x=[U(a),R]$,
\[
\partial_x(\chi_\pm^{-1}R\chi_\pm)
=\chi_\pm^{-1}(R_x-[U(a),R])\chi_\pm=0.
\]
另一方面, 式 \eqref{eq:Rbg} 和 \eqref{eq:diagE} 给出
\[
Y(z_a)^{-1}R_\infty Y(z_a)=\dfrac{\eta_0\lambda(z_a)}{a}\sg.
\]

\par\medskip\noindent\textbf{背景介质矩阵的特征值逐项核对.}\quad 
在共振原像有 $k=a$, 因而
\[
\lambda=a+\dfrac{q_0^2}{z_a},
\qquad z_a=a+\lambda.
\]
直接把 $R_\infty$ 与 $Y$ 相乘,
\begin{align*}
R_\infty Y
&=\begin{pmatrix}
\eta_0+\dfrac{\eta_0q_0^2}{az_a}&
-\dfrac{\ii\eta_0q_0}{z_a}+\dfrac{\ii\eta_0q_0}{a}\\
-\dfrac{\ii\eta_0q_0}{a}+\dfrac{\ii\eta_0q_0}{z_a}&
-\eta_0-\dfrac{\eta_0q_0^2}{az_a}
\end{pmatrix}\\
&=\dfrac{\eta_0\lambda}{a}
\begin{pmatrix}1&\dfrac{\ii q_0}{z_a}\\
-\dfrac{\ii q_0}{z_a}&-1\end{pmatrix}
=\dfrac{\eta_0\lambda}{a}Y\sg.
\end{align*}
例如右上元所用的等式是
\[
\dfrac{1}{a}-\dfrac{1}{z_a}
=\dfrac{z_a-a}{az_a}
=\dfrac{\lambda}{az_a}.
\]
左乘 $Y^{-1}$ 就得到原文的对角化公式.

$\lambda(z_a)=\pm\sqrt{a^2+q_0^2}$ 为非零实数, 空间指数及其逆均有界. 在各自归一化端取极限可得
\[
\chi_\pm^{-1}R\chi_\pm=\dfrac{\eta_0\lambda(z_a)}{a}\sg.
\]
再用 $\chi_+=\chi_-S_{\rm sp}$, 得到
\[
\dfrac{\eta_0\lambda(z_a)}{a}\sg
=S_{\rm sp}^{-1}\dfrac{\eta_0\lambda(z_a)}{a}\sg S_{\rm sp}.
\]

\par\medskip\noindent\textbf{从相同端态到反射为零的最后一步.}\quad 
在端点, 代入空间归一化可得
\begin{align*}
\chi_\pm^{-1}R\chi_\pm
&\longrightarrow e^{\ii\lambda x\sg}Y^{-1}R_\infty Y
 e^{-\ii\lambda x\sg}\\
&=\dfrac{\eta_0\lambda}{a}
 e^{\ii\lambda x\sg}\sg e^{-\ii\lambda x\sg}
=\dfrac{\eta_0\lambda}{a}\sg.
\end{align*}
指数与逆指数有界使趋零的端态误差仍然趋零.
两端计算出同一个对角矩阵, 再用散射关系可得
$S_{\rm sp}\sg=\sg S_{\rm sp}$.
把这个交换子逐元展开,
\[
S_{\rm sp}\sg-\sg S_{\rm sp}
=\begin{pmatrix}
0&-2(S_{\rm sp})_{12}\\2(S_{\rm sp})_{21}&0
\end{pmatrix}=0.
\]
所以 $(S_{\rm sp})_{12}=(S_{\rm sp})_{21}=0$.
再由
\[
1=\det S_{\rm sp}=(S_{\rm sp})_{11}(S_{\rm sp})_{22}
\]
可知两个分母都非零, 因而反射系数的零值确实有定义.
这里 $a\ne0,\eta_0\ne0$ 以及 $\lambda\ne0$ 都在除法时使用了.

因 $\dfrac{\eta_0\lambda(z_a)}{a}\ne0$, 两个非对角元必须为零. $\det S_{\rm sp}=1$ 保证对角元非零, 从而两反射系数均为零.
\end{proof}
